\section{Introduction}
\label{sec:intro}

Large language models (LLMs) have rapidly become the reasoning core of autonomous agents that write and execute code to solve data-science tasks~\citep{maojun2025survey,rahman2025llmbased}. Recent benchmarks such as DSBench~\citep{jing2024dsbench}, DataSciBench~\citep{zhang2026datascibench}, and DSAEval~\citep{sun2026dsaeval} report that even the strongest agents solve only 34--88\% of data-analysis tasks, leaving a large residue of failures. Yet these benchmarks report a single aggregate number---success rate---that conceals the \emph{structure} of failure: an agent that crashes on every hard task and an agent that silently returns plausible-but-wrong answers on the same tasks are indistinguishable under aggregate accuracy. This opacity matters because the two failure modes demand fundamentally different fixes: loud crashes call for better code generation, while silent wrong answers call for verification and self-correction mechanisms.

The problem is acute because silent failures are invisible to the agent itself. When generated code raises an exception, the agent observes the traceback and can retry; when the code executes successfully but produces a wrong answer, the agent has no signal that anything went wrong. Recent work has begun to recognise this distinction: \citet{mehta2026confident} report ``confident and wrong'' semantic failures in coding agents, \citet{wu2026when} document silent failures in production agent runtimes, and \citet{liu2026silent} frame silent failure as an entropy-driven inevitability. However, these works either study coding agents rather than data-science agents, or operate on production logs rather than controlled benchmarks, and none provides a \emph{diagnostic taxonomy} that localises where in the agent loop a silent failure originates or how far it propagates.

We address this gap with \textbf{AgentFail}, a diagnostic framework that decomposes every agent failure into four stages---\emph{Planning}, \emph{Tool-Use}, \emph{Execution}, and \emph{Interpretation}---and operationalises the silent/loud distinction at the taxonomy level. The Execution/Interpretation split is the key: a loud failure (Execution) is immediately visible to the agent and triggerable for retry, whereas a silent failure (Interpretation) is not. To move beyond observation and into \emph{attribution}, we introduce \emph{counterfactual causal replay}: given a failed trace, we replace the originating step with the ground-truth action and re-execute; if the outcome flips to correct, we have counterfactual evidence that the step caused the failure. This is strictly more informative than the observability-only tools criticised by \citet{shah2026causal}, which log what happened but cannot answer ``which step caused it?''

We evaluate five frontier LLMs (GPT-4o, GPT-4o-mini, DeepSeek-V3, DeepSeek-R1, Qwen3-max) across 4875 runs on three task suites of increasing realism: 80 controlled trap tasks designed to elicit specific failure modes, 23 real UCI datasets with analytical questions, and 200 real DSBench financial-analysis tasks. Our analysis surfaces four findings that aggregate accuracy cannot reveal:

\begin{enumerate}[leftmargin=*,itemsep=2pt]
\item \textbf{Failure bifurcation.} On solvable tasks, strong models (GPT-4o, DeepSeek-V3) exhibit a 100\% silent-failure rate---every failure is a wrong answer, never a crash---while weaker models fail loudly with execution errors; on very hard tasks (DSBench, $\sim$1\% success), all models fail loudly. Silent failure is thus the dominant failure mode \emph{precisely when agents are capable enough to run code but not capable enough to reason correctly about the result}.
\item \textbf{Zero self-correction.} The recovery rate---the proportion of detected failures that the agent subsequently corrects---is \emph{0\%} across all five models and all 4875 runs. Agents never recover from silent failures without external intervention.
\item \textbf{Causal attribution.} Counterfactual replay successfully attributes 79--88\% of failures to a specific originating step, providing actionable root-cause information that aggregate metrics cannot.
\item \textbf{Verifier boundary conditions.} A failure-aware verifier improves only the weakest model (Qwen3-max: +24.9\%, $p<0.001$, Cohen's $d=-0.72$) and is neutral or harmful on stronger models, indicating that silent failures resist simple automatic correction and motivating the need for more sophisticated verification.
\end{enumerate}

Our contributions are: (i) a four-stage failure taxonomy with an explicit silent/loud distinction; (ii) counterfactual causal replay for failure attribution; (iii) three complementary diagnostic metrics---silent-failure rate, propagation depth, and token-per-success---that expose structural failure patterns; (iv) the empirical discovery of failure bifurcation and zero self-correction across five frontier models; and (v) the public release of 4875 annotated traces and all framework code.\footnote{Anonymised repository: \url{https://github.com/llmnjust-afk/KDD2027_Work2}}

The remainder of the paper is organised as follows. \Cref{sec:related} reviews related work; \Cref{sec:method} presents the AgentFail framework; \Cref{sec:exp} describes the experimental setup and results; \Cref{sec:disc} discusses implications and limitations; and \Cref{sec:conc} concludes.
